% New navigation, summaries, citation context, and the compact control loop
% are isolated here. Extracted research text remains in the parent manuscript.

% BEGIN allocation-context
\subsection{Resource Allocation and Scheduling}

Resource allocation offers relevant control designs, but acts on different
objects. Chen et al. combine predicted Q-values with feedback control to
allocate cloud-service resources \cite{Chen2022RLFeedback}. Their
workload-window approach accounts for current and future demand
\cite{Chen2023WorkloadWindows}. Actor--critic allocation addresses dynamic
resource assignment in datacenters \cite{Chen2022ActorCritic}. Caviglione
et al. instead use DRL to choose among VM-placement heuristics with several
objectives \cite{Caviglione2021VMPlacement}. This is relevant to separating
action selection from execution. Here, the executed action changes a
subscription binding or a persistent connection, rather than VM placement.

Task schedulers also combine resource pressure with service objectives.
Ding et al. separate cloud-level task dispatch from per-server Q-learning
scheduling \cite{Ding2020QLearningCloud}. Chen et al. coordinate heterogeneous
workflows subject to task-continuity constraints
\cite{Chen2023CollaborativeScheduling}. Li et al. use DRL for energy-aware
Spark executor allocation under execution-time constraints
\cite{Li2023SparkDRL}. Cheng et al. study preemptive scheduling to reduce
cloud-job cost \cite{Cheng2023PreemptiveCloud}, while Liu et al. include
inter-region bandwidth and transfer cost in data-intensive task placement
\cite{Liu2026GeoDRL}. The review by Sanjalawe et al. places these scheduling
objectives within the wider AI-based cloud literature
\cite{Sanjalawe2025AIJobScheduling}. Our comparison concerns a different
execution requirement: preserving the evaluated transport endpoint while
transferring delivery ownership between existing client sessions.

Edge offloading provides related examples of constrained action selection.
Hortelano et al. classify RL approaches by their architectures, objectives,
and changing operating conditions \cite{Hortelano2023RLOffloading}. Huang
et al. learn binary offloading choices for wireless-powered devices
\cite{Huang2020Offloading}. Jayanetti et al. separate edge/cloud-tier selection
from node selection for precedence-constrained tasks
\cite{Jayanetti2022EdgeCloudScheduling}. These studies help distinguish the
decision space from the physical operation it controls. In this paper,
bounded templates select subscription redistribution or retained-probe
promotion, and receipt checks govern the resulting handoff.

% END allocation-context

% BEGIN adaptation-context
Monitoring and policy adaptation also need separate treatment. DEMon uses
decentralized monitoring with adaptive parameters to balance information
quality and resource use \cite{Ilager2026DEMon}. Our observation boundary is
narrower: each client sees its active endpoints and one retained probe.
Gheibi and Weyns study drift in the quality of available adaptation options
\cite{Gheibi2024AdaptationDrift}. This supports distinguishing a stored policy
from the conditions under which it was evaluated. Casimiro et al. use
probabilistic model checking to select adaptations of ML-based systems
\cite{Casimiro2024ModelChecking}. Our execution guards check concrete handoff
conditions; they are not a model-checking guarantee for the learned policy.

% END adaptation-context

% BEGIN evaluation-context
RL evaluation also separates adaptation from regression. Biagiola and Tonella
test how RL systems adapt to changed environments while retaining prior
capabilities \cite{Biagiola2022Plasticity}. Patterson et al. discuss performance
variation, baseline construction, hyperparameter selection, and experimenter
bias in empirical RL \cite{Patterson2024EmpiricalDesign}. These concerns motivate
reporting the training, frozen-policy comparison, and live handoff evidence
separately. A successful handoff and a higher comparative reward answer
different evaluation questions.

% END evaluation-context

% BEGIN related
The broader cloud-scheduling and edge-computing context is retained in
Section~\ref{app:extended-literature} of the Online Supplementary Material.
% END related

% BEGIN system
\section{System Model and Receipt-Gated Handoff}
Figure~\ref{fig:control_architecture} shows the shared architecture. Event
producers send messages through ingest services to durable subscription queues.
Subscriber Pods deliver them to clients over persistent connections. The
transport load balancer assigns new sockets; the client selects whether to
promote its retained probe. Additional component descriptions are provided in
the Online Supplementary Material.
% END system

% BEGIN model
Detailed reference-scale calibration, the global evaluation objective, and
step-by-step proofs of the subscription-density and backlog-variance
properties are provided in Section~\ref{app:extended-mathematics} of the Online
Supplementary Material. The full notation reference is in
Section~\ref{app:notation}.
% END model

% BEGIN contract
Here $\Theta_N$ is the configured maximum desired subscription count per
connection, not the network parameter vector $\theta$.
The evaluated legacy checkpoint uses a compatible action profile; both
scaling actions are disabled in the loaded cluster trial. The later
cap-restoring scale-up contract is distinct. Table~\ref{tab:action-contract-map}
retains these version boundaries and their evaluation scope.
% END contract

% BEGIN algorithm
\subsection{Shared Control Loop}
Algorithm~\ref{alg:compact-control} summarizes the policy-independent runtime.
The full sizing and rebalancing executors and the DQL training loop are
retained in Section~\ref{app:runtime-details} of the Online Supplementary Material.

\begin{algorithm}[!htbp]
\small
\caption{Probe-Guided Selection and Guarded Execution}
\label{alg:compact-control}
\SetKwInOut{Input}{Input}
\Input{Active connections, retained probe, declared policy and action profile}
\While{client is active}{
 Read active/probe telemetry and client queue/ACK counters\;
 Construct observation and freshness, capacity, cooldown and handoff masks\;
 Select an admitted action using the declared greedy or DQL policy\;
 Recheck guards before applying the selected action\;
 \eIf{migration is selected and admitted}{
  Prepare standby bindings on the same probe; await receipts\;
  Quiesce the source; wait for fresh zero-in-flight telemetry\;
  Activate target bindings; detach source bindings; retire the source socket\;
  Retain unselected sessions and replenish the empty probe\;
 }{
  Execute the admitted bounded rebalancing/scaling action, or stay\;
 }
 Record the applied action, observed reward and guard outcome\;
 Update learning state only in a declared training run with a valid transition\;
}
\end{algorithm}
% END algorithm

% BEGIN runtime
The exact network input order, normalization constants, and deterministic
comparator thresholds are retained in Section~\ref{app:feature-order} of the
Online Supplementary Material. Checkpoint persistence and full execution
details are in Section~\ref{app:runtime-details}.
% END runtime

% BEGIN timing
The paired handoff/drain distributions and their statistical interpretation
are retained in Section~\ref{app:cloud-timing} of the Online Supplementary Material.
% END timing

% BEGIN ttl
\subsection{TTL Reconnection Comparison}
An isolated matched-preload pair used 5,000 subscriptions and 30,000 delivery
attempts per arm. Native TTL produced twelve server-matched interruptions and
26,831 confirmed ACKs (89.44\%), leaving 3,169 unresolved receipts. The
TTL-disabled DQL arm confirmed all 30,000 ACKs without interruptions, but its
telemetry guards held every decision at \emph{stay}. A 200-second minimum
telemetry sweep exceeded the 180-second freshness deadline. This pair measures
expiry-related interruption and receipt completion, not a learned-migration
throughput advantage. Unresolved receipts are not classified as event loss.

Detailed TTL implementation paths, recovery endpoints, the two-burst trial,
and matched-preload diagnostics are provided in
Section~\ref{sec:ttl-reconnection} of the Online Supplementary Material.
The separate 2,000-subscription calibration and startup-under-load trial is
retained in Section~\ref{sec:mumbai-startup}.
% END ttl

% BEGIN scaleout
Capacity-activation timelines, seed-level violin plots, and independent
ad-hoc timing controls are retained in Section~\ref{app:scaleout-details} of
the Online Supplementary Material. Those controls do not replace the
pressure-aware guarded-greedy comparator above.
% END scaleout

% BEGIN holdout-heading
\subsection{Independent Multi-Objective Holdout}
The independent holdout compares the frozen DQL policy with the threshold-based
Heuristic and Static policies from Table~\ref{tab:policy-roles}. Full study
designs, training traces, distributions, and action-contract diagnostics are
retained in Section~\ref{app:holdout-details} of the Online Supplementary Material.
% END holdout-heading

% BEGIN ablation
\subsection{Adverse-Probe Ablation}
The separate 30-seed, 1,200-subscription ablation compares the same frozen DQL
policy with and without one forced promotion to the highest-pressure probe
Pod. Guarded DQL reduces normalized backlog AUC from 0.68938 to 0.66046:
a paired difference of $-0.028918$ (95\% CI $[-0.039444,-0.019073]$), or 4.19\%.
The guarded heuristic retains the lower AUC of 0.59104. This isolates the
benefit of guarding promotion, not universal superiority of learned selection.
The complete five-policy comparison, methods, and separate five-seed
20,000-subscription confirmation are retained in
Section~\ref{sec:adverse-probe-ablation} of the Online Supplementary Material.
% END ablation

% BEGIN repository

The public research repository is available at
\url{https://github.com/manish94u/rebalancePersistance}.
It provides the manuscript sources, standalone simulator, redacted evidence,
frozen checkpoint, and step-by-step replication instructions.
% END repository

% BEGIN conclusion
The Oracle Cloud experiments establish live handoff execution, subscription
retention, and observed delivery behavior. The comparative throughput and
backlog-recovery gains are simulator measurements under the stated service
assumptions. These results support two evaluated policy options above the same
probe-guided handoff mechanism.
% END conclusion

% BEGIN math-heading
\section{Extended Mathematical Details}
\label{app:extended-mathematics}
The core queue, pressure, imbalance, and reward definitions appear in the main
paper. This section retains their reference-scale calibration, the offline
objective, and the full proofs.
\subsection{Variance Reference Scales}
% END math-heading

% BEGIN runtime-heading
\section{Detailed Runtime Algorithms and Checkpoint Persistence}
\label{app:runtime-details}
The main paper gives the shared control-loop summary. The following material
retains the action-contract details, checkpoint behavior, sizing rules,
executors, and full DQL algorithm.
\subsection{Action Contract}
% END runtime-heading

% BEGIN literature-heading
\section{Extended Cloud and Edge Literature}
\label{app:extended-literature}
% END literature-heading

% BEGIN cloud-heading
\section{Additional 20,000-Subscription Handoff Diagnostics}
\label{app:cloud-timing}
The deployment setup and aggregate results appear in the main paper. The
following paired plot retains the within-run timing detail.
% END cloud-heading

% BEGIN scaleout-heading
\section{Additional Scale-Out and Timing-Control Results}
\label{app:scaleout-details}
The workload, service model, primary comparison, and backlog trajectories
appear in the main paper. This section retains the capacity-activation plots,
distributions, and exploratory timing controls for the same experiment.
% END scaleout-heading

% BEGIN holdout-detail-heading
\section{Full Supporting Simulator Studies}
\label{app:holdout-details}
Study designs and explanatory passages are retained with their original
results. Selected definitions and results also appear in the main paper so
that its core comparison can be read independently.
% END holdout-detail-heading

% BEGIN discussion-heading
\section{Extended Interpretation}
The complete interpretation from the reference manuscript is retained below.
% END discussion-heading

% BEGIN supp-front
\settopmatter{printacmref=false}
\maketitle
\raggedbottom
This Online Supplementary Material accompanies the main manuscript
\emph{Probe-Guided Rebalancing of Persistent Subscription Connections with
Greedy and Deep Q-Learning Control}. It retains the detailed mathematical,
implementation, and experimental material relocated from the reference
manuscript. Reported measurements, comparisons, and figures are unchanged.
References explicitly marked ``of the main paper'' point to the separate
main manuscript. Figures, tables, equations, and algorithms in this document
use an S prefix.
\appendix
\renewcommand{\thefigure}{S\arabic{figure}}
\renewcommand{\thetable}{S\arabic{table}}
\renewcommand{\theequation}{S\arabic{equation}}
\renewcommand{\thealgocf}{S\arabic{algocf}}
\setcounter{tocdepth}{1}
\tableofcontents
\clearpage
% END supp-front
